% Subproblem 4 (EASY): Convert Xi >= 0 into I < 0
%
% CONTEXT: Part of the T_{2g} Hessian proof.
% Working directory: /data/projects/3eacb340-027e-4b71-8de4-67574ba54ea2/QBLYHzbdEjWeyg/proof/
%
% GIVEN:
%   - rho(theta, psi) = rho^+(theta,psi) + rho^-(theta,psi) where rho^pm = exp(+/-u)/I^pm,
%     and u = 2*pi * (G_A + G_B) is the log-norm function. rho > 0 almost everywhere.
%   - Xi(theta, psi) = (1/2)*D(theta,psi)^2 + (1/2)*D(theta,psi+pi/2)^2 - C(theta,psi)^2
%     is locally integrable on S^2, nonneg for (theta, psi) in (0,pi) x (0,pi/2) away
%     from the cube vertices (proved in Subproblem 3), and not identically zero on any
%     open set.
%   - By Subproblem 2:
%       I = integral_{S^2} F_A * F_A(psi+pi) * d(mu^+ + mu^-)
%         = - integral_0^pi integral_0^{pi/2} Xi(theta,psi) * rho(theta,psi) * sin(theta) d psi d theta.
%
% STATEMENT TO PROVE:
%   I < 0.
%
% PROOF:
%   Since Xi >= 0 a.e. in (0,pi) x (0,pi/2), and rho * sin(theta) is a positive measure
%   on that domain (rho > 0 a.e.), the integral integral Xi * rho * sin(theta) dpsi dtheta >= 0.
%   Since Xi is not identically zero (it is strictly positive on some open set of positive
%   measure), and rho > 0 on a set of full measure, the integral is strictly positive.
%   Therefore I = -(positive number) < 0.
%
% DELIVERABLE: A short, clean LaTeX proof (2-5 sentences) of I < 0, citing Subproblems 2 and 3.
% Output should be a self-contained LaTeX fragment (no \documentclass).
\end{document}
